\chapter{Related Works}
\label{chap:final-related-works}

\section{Metallic Artifacts in\acrshort{ct}}

The traditional methods of artifact reduction in\acrshort{ct} are predominantly based on physical and mathematical principles associated with image formation. Barrett and Keat\cite{barrettArtifactsCTRecognition2004} provide a detailed description of the origin of the main artifacts in\acrshort{ct}, including beam hardening, photon starvation, scattering, and geometric effects, as well as classical strategies for their mitigation.

Reconstruction techniques such as\acrshort{fbp} were widely used due to their computational efficiency, but they are highly sensitive to noise and metallic artifacts. More advanced methods, such as\acrfull{ir} and\acrshort{mbir}, incorporate statistical models of the acquisition system and noise, allowing for partial artifact reduction. However, these approaches often introduce artificial textures and have limitations in complex clinical scenarios \cite{sahuTransformingCTImaging2025}.

Recent revisions indicate that commercial algorithms for\acrshort{mar}, such as\acrfull{omar},\acrfull{imar}, and\acrfull{semar}, although effective in specific situations, rely on proprietary parameters and exhibit inconsistent performance in cases of dense or multiple implants \cite{kleberAdvancementsSupervisedDeep2024, sellesAdvancesMetalArtifact2024}. These limitations have driven the investigation of deep learning-based approaches. 

\section{Artifact Reduction based on Deep Learning}

The First applications of\acrfull{dl} para\acrshort{mar} focused on the image domain, exploring convolutional networks to map degraded images to corrected versions. Gjesteby et al.~\cite{gjestebyDeepLearningMethods2017} investigated metal artifact reduction in the image domain based on\acrshort{cnn} and demonstrated significant quantitative improvements compared to conventional techniques.


Zhang e Yu~\cite{zhangConvolutionalNeuralNetwork2018} also purposed a supervised model based on\acrfull{cnn} for metal artifact reduction in\acrshort{ct}, achieving consistent gains in\acrshort{psnr} and\acrfull{rmse}. These works helped consolidate the use of convolutional networks for\acrshort{mar} in the image domain. In particular, Zhang and Yu used a shallow fully convolutional\acrshort{cnn} for fusing information from multiple pre-corrected images. Subsequent works increasingly adopted encoder-decoder architectures, especially those based on\acrfull{unet}, for correction/restoration tasks. \acrshort{unet}-type networks and their variants are capable of capturing global contextual information through the contraction path, and recovering fine spatial details through the expansion path and \textit{Skip Connections}\cite{ronnebergerUNetConvolutionalNetworks2015}.

Recent Sistematic reviews have highlighted these advances, emphasizing that most supervised methods in the image domain exhibit competitive performance, especially when trained with high-quality paired data \cite{kleberAdvancementsSupervisedDeep2024,sahuTransformingCTImaging2025}.

In addiction to image domain, some studies explored the sinogram domain, with linear interpolation and\acrfull{nmar} methods. These approaches operate directly on the projection data and act directly on the physical origin of metallic artifacts\cite{meyerNormalizedMetalArtifact2009}. However, access to raw projection data is often restricted in commercial\acrshort{ct} systems, limiting the applicability of such approaches in clinical settings \cite{gjestebyDeepLearningMethods2017}.

Beyond that, hybrid models have been purposed to leverage the advatages of each approach. The method purposed by Zhang and Yu~\cite{zhangConvolutionalNeuralNetwork2018} seeks to merge\acrshort{cnn} in the image domain and classical\acrshort{mar} in the sinogram domain. In the adopted approach, the initial refinement occurs in the image domain, but the complete method uses the image generated by the\acrshort{cnn} as a prior for replacement of affected regions in the sinogram. This characterizes a hybrid approach, but still does not employ\acrshort{cnn} in both domains. More recent models, such as DuDoNet\cite{linDuDoNetDualDomain2019}, integrate deep networks in both domains, demonstrating high performance and characterizing what we call Dual Domain. However, these methods present greater computational complexity and dependence on access to raw sinogram, which motivates the continuous exploration of robust solutions in the image domain.

To address the scarcity of anatomically paired data with metal artifact/metal artifact-free, self-supervised learning-based approaches have emerged. Yu et al.~\cite{yuMetalArtifactReduction2021} proposed a method capable of reducing metal artifacts without the need for fully clean reference images, expanding the clinical applicability of these techniques.

Others works explored the use of\acrfull{gan} for\acrshort{mar}. Huang et al.~\cite{huangMultimodalFeaturefusionCT2022} used multimodal fusion and\acrshortpl{gan} to improve texture preservation and reduce excessive smoothing. Despite good visual results,\acrshort{gan}-based can introduce training instability and a higher risk of structural hallucination\cite{yiGenerativeAdversarialNetwork2019}, raising clinical concerns. 


\section{Reproducible Evaluation and Generalization}

The evaluation of\acrshort{dl} models for\acrshort{ct} is typically performed through quantitative metrics such as\acrshort{psnr},\acrshort{rmse}, and structural similarity-based indices. Traditional error-based metrics, such as\acrshort{rmse} and\acrshort{psnr}, quantify pixel-wise discrepancies but do not adequately capture structural information relevant to visual perception and clinical interpretation. In this context, Wang et al.~\cite{zhouwangImageQualityAssessment2004} introduced the\acrshort{ssim} as a perceptual metric that explicitly models luminance, contrast, and structure. Subsequently, the\acrfull{ms-ssim} was proposed to incorporate scale-dependent structural variations, extending\acrshort{ssim} to multiple resolution levels \cite{wangMultiscaleStructuralSimilarity2003a}. Hor{\'e} and Ziou~\cite{horeImageQualityMetrics2010} demonstrated that both\acrshort{psnr} and\acrshort{ssim} have limitations depending on the type of distortion considered, suggesting that the choice of metric may depend on the application context. Given that metallic artifacts manifest primarily as structural distortions affecting anatomical boundaries and long-range streak patterns, metrics sensitive to structural preservation, such as\acrshort{ssim} and\acrshort{ms-ssim}, are particularly suitable for evaluating the performance of\acrshort{mar} methods, as they explicitly consider structural components at different spatial scales~\cite{zhouwangImageQualityAssessment2004,wangMultiscaleStructuralSimilarity2003a}.  

Additionally, a critical challenge identified in the literature is the models generalization. Zech et al.~\cite{zechVariableGeneralizationPerformance2018} demonstrated that models with high performance on internal datasets can fail significantly when applied to external data, highlighting the impact of Domain Shift. Recent reviews reinforce the need for rigorous external and clinical validation to ensure safety, robustness and real applicability~\cite{willeminkPreparingMedicalImaging2020,kellyKeyChallengesDelivering2019}. 

The reviewed literature demonstrates significant advancements in artifact reduction in\acrshort{ct} through\acrshort{dl}, but also evidences persistent gaps: lack of experimental standardization, limited external validation, and risk of anatomical hallucination.